Redis节点内存中已经存在一整份数据时，RDB把数据编码为可重新加载的快照。查询缓存可以舍弃快照并重建，较大的投影也可以借助快照缩短冷启动时间。沿上一节的第二条假设路径，若Redis直接承担权威责任，恢复目标就要包含完整身份和所有尚未结束的工作：关闭会话、S绑定、原输入、P/C/T、原平台scope/\code{pay-C}、原结果、O/EV、消费交接回执与进度、退款占额和库存事务责任，都要能从相应保存域取得。只备份页面看见的\code{paid}状态，会遗漏付款未知时最需要恢复的关联~\cite{redis8102-persistence}。

\subsection{同步保存与后台保存}

\code{SAVE}在当前进程内遍历并写出数据，保存期间阻塞普通命令处理。\code{BGSAVE}通过fork创建子进程，子进程看到fork时点的地址空间，父进程继续接收请求\footnote{fork是操作系统创建子进程的接口。父子在逻辑上分别拥有地址空间，物理页可以在写入前共享；写时复制在写入共享页时安排独立页面。}。两者的差别首先在于谁承担遍历与编码，快照内容都需要对应一个可说明的数据状态。

设fork之前P1的投影是版本1，父进程在fork之后改成版本2。操作系统最初允许父子共享物理页，父进程要改写共享页时，通过写时复制为相应修改保留自己的页面。子进程继续读取旧页面，所以RDB得到fork时点的版本1。父进程后来新建的键不会因为保存任务尚未结束就自动进入这次快照。快照完成时间决定文件何时可用，数据时点由fork边界决定。

\CEighteenFigure{ch18-05}{BGSAVE的共同内存前缀与写时复制}{fig:18-rdb-cow}

后台保存使遍历和磁盘写入与父进程命令处理重叠，但fork、页表复制、写时复制和磁盘竞争都有成本。数据集较大、保存期间修改密集时，父进程的新增页面和子进程工作集会增加内存压力。把可用内存全部留给键值，缺少后台保存余量，可能使一次保存触发内存耗尽。节点容量因此同时计算数据、客户端与复制缓冲、后台任务和操作系统所需空间。

\subsection{从临时文件到完成文件}

RDB保存先写临时文件，写出必要头部、对象内容及结束校验信息。完成路径检查写入、flush、文件同步和关闭等结果，再将临时文件重命名到目标文件，并处理目录同步。这个顺序让原完成文件在新文件写到一半时继续保留。子进程中途退出，父进程根据退出状态处理失败，半份临时文件不会因此成为一份成功发布的新快照~\cite{redis8102-source}。

\CEighteenFigure{ch18-06}{新RDB完成后替换旧文件，失败分支保留旧完成文件}{fig:18-rdb-publish}

文件系统的原子重命名负责名称切换，持久性还依赖文件与目录同步、文件系统及存储设备实际提供的保证。把文件复制到其他机器时，复制工具也要先取得完整文件，再核对大小与摘要并发布完成标记。正在传输的一段字节，即使文件名已经叫\code{dump.rdb}，也要等完整性核验后才进入可用恢复清单。

应用若把完整分区编码成多个键，这次快照对Redis本地数据形成共同切面，但它只能保存fork时已存在的键。假如应用先保存S占位，尚未保存T就执行后台快照，RDB会忠实保存这个不完整的业务状态。业务更新先在命令执行边界内形成可恢复的完整受理结构，再让快照、AOF和复制承载它。第17章冻结图像进一步撤销源分区的写资格，用于跨组交接完整权威。

快照中的完整性可以按关系逐层核对。会话E1记录账户、动作、固定截止及关闭事实，$A_J$由完整S定位到原输入和P1/C1/T1，T1指向原平台scope/\code{pay-C1}并说明当前恢复阶段。若已有原付款证据，$R_J$指向同一C1/X1以及原结果，O1与EV1保存对应事件内容及发布责任。已经交接的EV1还关联原消费回执、履约来源与进度，快照据此恢复补发是否停止。退款对象再关联原收款和已成功、未确定的占额。核对这些关联指向同一业务历史，才决定图像能否接管。

快照完成与执行资格分别管理。A取得源B数据的一份RDB，源B同时继续处理T1，A若立即启动自己的付款工作者，就形成两个执行者。原平台键可以吸收同一操作的重发，资源写入、退款占额和新受理还有各自的资格约束。备份恢复先把目标放在隔离状态，核验图像与外部证据，再按控制面的接管规则开放相应动作。架构中分别登记数据副本和活动权威。

\subsection{校验、加载与过期时间}

加载器检查格式和版本，按记录类型重建键及对象，并在相应配置下核对文件校验和。校验和能发现相应字节损坏，文件来自哪个业务分区、属于哪个历史时点，则由恢复清单中的来源、快照时点、适用配置和摘要共同确定。

RDB保存键的过期信息，普通主节点恢复时可以跳过已经过期的键。复制加载或AOF的RDB前导段采用与后续流衔接的规则，过期处理要和主副本角色配合，保证主节点传播的修改能够形成同一条历史。应用把关键准入截止和关闭事实作为业务字段保存，TTL只管理承载键的生存时间~\cite{redis8102-source}。

键级过期与业务关闭具有不同恢复含义。缓存项已过期，重新读取当前权威即可。会话E1已关闭，关闭事实要持续保留在该业务采用的准入与身份保留规则中，恢复时据此拒绝重新开放E1。保存截止值帮助判断相应时间条件，已经持久发生的关闭事件另由关闭记录表达。恢复者还核对节点时钟是否满足部署假设，避免把错误时钟造成的TTL判断当作业务历史。

RDB文件损坏时，恢复者先保留原文件与获取记录，在副本上检查。能够解码部分对象，只说明这部分字节能够读取。设恢复出的数据含C1，却缺少它的原平台scope/\code{pay-C1}，工作者就失去了安全查询X1的依据。系统从其他完整副本或可验证图像恢复关联，或者冻结相关动作并对账，直到找回原平台操作的身份。

\subsection{旧快照与独立外部事实}

初态是某份RDB已经保存$A_J$中的P1/C1/T1，T1处于付款未知。随后平台按\code{pay-C1}产生成功交易X1，权威又保存原结果与O1/EV1。灾难恢复只能取得前面的RDB时，Redis恢复了旧数据时点，平台的X1继续存在。工作者从原C1/T1读到原scope/\code{pay-C1}，查询平台原操作，再把同一证据汇合到原购买，才有机会重建后面的本地结果。

\CEighteenFigure{ch18-07}{旧快照恢复本地责任，平台原操作保留在独立历史中}{fig:18-old-rdb-external}

若快照连C1的首次受理都没有包含，恢复问题更早发生：旧S是否已经发送到外部、哪些编号已经分配，都可能超出本地恢复点。此时沿用第14章的灾难恢复规则，冻结存在歧义的受理范围，从跨故障域证据恢复旧 S 的原绑定、原 C 和原结果；无法恢复的范围通过持久关闭规则拒绝旧提交。对经核验可以建立的真正新意图，另从不会重用的来源取得编号。更高代次和新编号只解决编号复用，旧 S 的历史由前一道恢复或关闭规则处理。普通副本切换下连续的计数器，在旧快照恢复后可能回退，分区前缀只能区分分区，不能补回已经消失的分配历史。

回到购买系统采用的缓存路径，RDB落后只影响投影的恢复起点。恢复者先清空不可信投影，再从当前权威取得相应版本；携带最低版本的查询在新投影达标前回源。
\FloatBarrier
